
\subsection{Overfitting Controls in Skill Synthesis}
\label{app:synthesis-overfit}

Synthesizing effective defense skills requires balancing two opposing
risks: over-specialization to specific attack variants and
over-generalization that triggers false positives on benign
traffic. System~II enforces this balance by decoupling the validation
constraints applied to candidate skills from the information exposed
to the generator.

Here the verifier serves as the strict quality gate. Specifically, a
candidate skill is accepted only if it matches all attack trigger
samples while firing on no more than one of benign samples. To prevent
Denial-of-Service via regex exploitation, checks execute in an
isolated child process subject to a 200~milliseconds timeout alongside
static regex analysis. Crucially, the generator operates across an
information barrier: it receives only a structured abstract profile of
the attack (e.g., length, sample count, keyword matches, Shannon
entropy, and encoding flags), the current skill graph metadata, and
failure feedback from prior rounds, rather than the raw trigger
text. To broaden coverage without compromising validation rigor, a
small cluster of nearest-neighbor attack samples is provided as
contextual guidance, though these are excluded from the formal pass
predicate.
When an entire synthesis round fails, the engine transitions from
fresh generation to targeted mutation on the best candidate marked for
revision. Once accepted, candidates lacking concrete payloads undergo
post-acceptance gating: they enter a shadow observation mode and
remain non-blocking until completing a seven-day or 50-scan clean
window with zero false positives and at least one confirmed
hit. Finally, as a fail-safe against overly broad signatures that
bypass earlier checks, such as predicates matching generic environment
content, a pre-deployment filter prunes loose rules before active rollout.

